% =====================================================================
%  12 -- THE BASELINE YOU ARE TRYING TO BEAT  (source node S16, Part VI).
%
%  Source: tier3.md lines 4612-4909, in full and in order --
%          Part VI header (4612-4615), VI.1 (4617-4673), VI.2 (4674-4728),
%          VI.3 (4729-4757), VI.4 (4758-4790), VI.5 (4791-4837),
%          VI.5b (4838-4894), VI.6 self-check (4895-4909, six questions ->
%          sec:selfcheck-baseline).
%
%  Rendering notes:
%    - the MLP arrow diagram (source 4627-4630) is the assigned TikZ
%      figure \label{fig:nnnarch}; the two python fences (log-softmax,
%      loss) are srclisting verbatim;
%    - the novelty table (4779-4786) is set as tabularx
%      {L{2.1cm} L{2.6cm} L{4.6cm} Y} \small.
%  Re-homings: none.  No cross-reference sentence had to be reworded
%  (Parts III, IV, V and sections 0.3.6, 0.6.5, I.5, I.6 all precede
%  this part in the report, as they do in the source).
%  Tier 2 / Tier 0 pointers stay as plain text per the brief.
% =====================================================================

\section{The baseline you are trying to beat}
\label{part:baseline}

\textbf{Know your competitor better than they do.} This node is the NNN as an
object, what it achieves, where it does not, and what ``beating it'' has to
mean.

% =====================================================================
\subsection{What the NNN actually is}
\label{sec:nnnis}

Grichener, Renzo, Kerzendorf et al. 2025, \emph{Nuclear Neural Networks:
Emulating Late Burning Stages in Core-collapse Supernova Progenitors} (ApJS
279, 49; \arxiv{2503.00115}; data and models at Zenodo 10.5281/zenodo.14873443)
\citep{grichener2025}. Read from the shipped code
(\code{repro/nnn/patched/NNNfunctionsCompPlusEps.py}) and checked against the
paper (\S 2.3):

\textbf{Architecture.} A plain feed-forward MLP:

\begin{figure}[H]
\centering
\begin{tikzpicture}[
  node distance=5mm and 9mm,
  box/.style={draw, rounded corners=1.5pt, thick, font=\small\ttfamily,
              minimum height=7mm, inner xsep=5pt, align=center},
  act/.style={font=\footnotesize\ttfamily},
  arr/.style={-{Latex[length=2.2mm]}, thick}]
  \node[box] (in)  {input\\(n\_iso + 2)};
  \node[box, right=of in]  (l1) {Linear\\1024};
  \node[box, right=of l1]  (l2) {Linear\\2048};
  \node[box, right=of l2, draw=none, font=\large] (dots) {$\cdots$};
  \node[box, right=of dots] (l3) {Linear\\2048};
  \node[box, right=of l3]  (out) {Linear\\(n\_iso + 2)};
  \draw[arr] (in) -- (l1);
  \draw[arr] (l1) -- node[act, above] {ReLU} (l2);
  \draw[arr] (l2) -- node[act, above] {ReLU} (dots);
  \draw[arr] (dots) -- (l3);
  \draw[arr] (l3) -- node[act, above] {ReLU} (out);
\end{tikzpicture}
\caption{The NNN architecture as shipped: a plain feed-forward MLP,
  input ($n_{\mathrm{iso}} + 2$) $\to$ Linear 1024 $\to$ ReLU $\to$ Linear
  2048 $\to$ ReLU $\to$ \ldots\ $\to$ Linear 2048 $\to$ ReLU $\to$ Linear
  ($n_{\mathrm{iso}} + 2$).}
\label{fig:nnnarch}
\end{figure}

Twelve (\msn{80}) or nine (\msn{151}) Linear layers: input $\to$ 1024 $\to$
2048 ($\times 10$ or $\times 7$) $\to$ out (\citealt[\S 2.3, Fig.~2]{grichener2025};
\code{docs/zenodo-inventory.md}). The class constructor's defaults --- four
layers of 128/256/256 with \code{neuronsNumFactor = 1, layersNum = 4} --- are
\emph{not} the shipped configuration, which is loaded with factor 8 and 12/9
layers (\code{runNNNsOnTestMesa80.py:491}); an earlier version of this section
quoted the defaults (corrected in the 2026-08-18 audit). Input = the initial
mass fractions plus $\log T$ and $\log\rho$ (dims 82 / 153, confirmed against
the checkpoints \resultsref{2026-07-08}). Output = the final mass fractions
plus \enuc{} and \epsnu{} (the energies normalised by $10^{16}$).

\textbf{The output transform is the interesting part.} The isotope block
passes through a \textbf{log-softmax}:

\begin{srclisting}
first_80_log_softmax = torch.log_softmax(x[:, :isotopesNum], dim=1)
x = torch.cat([first_80_log_softmax, x[:, isotopesNum:]], dim=1)
\end{srclisting}

So the predicted mass fractions satisfy $\sum\exp(\mathrm{out}_i) = 1$
\textbf{exactly, by construction} (in float32; \citealt[\S 2.3]{grichener2025}:
``to ensure the abundances sum to one''). The NNN \emph{does} conserve mass ---
through a normalisation, not through stoichiometry.

\textbf{The loss} is L1 in log space on the isotopes plus L1 in linear space on
the two energy outputs \citep[\S 2.4]{grichener2025}:

\begin{srclisting}
isotopes_loss = F.l1_loss(nnComp_isotopes, torch.log(bbqComp_isotopes))
eps_loss      = F.l1_loss(nn_eps, bbq_eps)
total_loss    = isotopes_loss + eps_loss
\end{srclisting}

\textbf{Nine models per network.} One MLP per timestep --- the dt grid is not a
model input, it is a model \emph{index} (\citealt[\S 2.2]{grichener2025}: ``a
separate neural network for every timestep''; the paper's burner interpolates
between the two nearest trained dt, \S 4). Eighteen production models ship
(7.3~GB; the 23.8~GB \code{trained\_NNN\_models/} tree also holds a 36-model
layer scan, each checkpoint carrying Adam optimizer state).

\textbf{Training data.} The $2^{20}$ Sobol grid, 1,041,400 usable states per
network per dt (\S\ref{sec:shippedpipeline}), cast to float32 by the loader.

% =====================================================================
\subsection{Reading that design critically}
\label{sec:nnncritical}

Four observations, each of which is a design axis this project takes
differently.

\textbf{(i) Mass is conserved; charge is not.} The softmax fixes
$\sum X_i = 1$ exactly and says nothing about $\sum Z_i Y_i$. \textbf{\Ye{} is
entirely unconstrained}, which is why the paper reports a \Ye{} \emph{loss} as
a separate metric rather than an invariant (\citealt[\S 3]{grichener2025},
their $\delta\Ye = \sum_i (Z_i/A_i)|\Delta X_i|$ --- an upper bound on the
signed error), and why the reproduced relative \Ye{} error is 0.33--0.54\%
(\msn{80}) / 0.55--0.73\% (\msn{151}) \resultsref{2026-07-08} (the paper
quotes 0.4--0.75\%).

Put that number in the project's own units. $\Ye \approx 0.47$, so 0.5\%
relative is \textbf{$|\dYe| \approx 2.4 \times 10^{-3}$ per step} (the measured
band gives 1.6--3.4 $\times 10^{-3}$). The project's per-step gate is
$3 \times 10^{-6}$. \textbf{The baseline is ${\sim}800\times$
(520--1140$\times$) outside the gate this project sets itself.} That is a
startling gap, and it is the single most important number for calibrating
ambition. Two honest readings coexist:

\begin{itemize}
\item the gate may be over-specified (\S\ref{sec:missingderivative} /
  \reg{T-01} --- the derivative that would tell you is missing);
\item or the NNN's \Ye{} accuracy is genuinely inadequate for the science,
  which is exactly the claim this project is built on.
\end{itemize}

Either way, \textbf{this comparison should be in the project's own documents
and is not.} Registered as \reg{T-07}.

\textbf{(ii) Normalisation is not conservation.} Softmax rescales; it does not
respect stoichiometry. Tier~2 \S 0.5 makes the general argument: a leak is
species-specific, a renormalisation is uniform, so renormalising after a leak
silently redistributes abundance between species. The NNN never ``leaks'' mass
because it never computes a change --- it predicts the final state directly ---
but the same criticism applies in the form: \emph{there is no mechanism forcing
the predicted state to be reachable from the initial one.} Charge, lepton
number, and the reachable-manifold structure (Tier~2 \S 0.4.3) are all
unconstrained.

\textbf{(iii) The loss is in log space, which is the right choice for dynamic
range and the wrong space for constraints.} Mass fractions span 1e-15 to 1; an
L1 loss in linear space would be dominated entirely by the top few species. Log
space fixes that and is a genuinely good decision. But it means the loss
weights a relative error on a $10^{-14}$ species the same as on a 0.5
species --- consistent with the paper's per-isotope error band of $10^{-4}$
\ldots\ $10^{-1}$ with light isotopes \emph{best} (\citealt[\S 3]{grichener2025}
attribute the ordering to abundance: absolute errors grow with $X$) --- and it
means no linear constraint can be imposed in the space where the loss lives.
This is the same structural point as invariant \#4, arriving from the loss
side.

\textbf{(iv) Nine models is a strong statement about the hard part.} Training
one model per dt admits that the $\Delta t$ dependence is the difficult axis.
Tier~4's S19 takes the opposite bet --- one model over eight decades, with
$\Delta t$ as conditioning --- and \foreignreg{2}{R-10} argues the conditioning
variable should be $\Delta t/\tau$ rather than $\log\Delta t$ precisely because
\S\ref{sec:realtrajectory} measured that \emph{every} label is a full
relaxation. Nine models is the safe design; one conditioned model is the better
one if it works, and it is $9\times$ cheaper to store and deploy.

% =====================================================================
\subsection{What the NNN achieves}
\label{sec:nnnachieves}

Reproduced exactly by this project (\code{repro/nnn/}), which is worth its own
paragraph. \resultsref{2026-07-08}:

\begingroup\small
\begin{tabularx}{\textwidth}{@{}L{5.0cm} Y L{3.6cm}@{}}
\toprule
\textbf{claim} & \textbf{reproduced} & \textbf{published} \\
\midrule
all 18 (net, dt) $\times$ 6 metrics vs shipped result CSVs & \textbf{ratio
  1.000 in all 108 comparisons} & --- \\
independent state\_dict-walk loader vs upstream class, 128 samples &
  \textbf{max deviation 0.0} (gate $\le 10^{-6}$) & --- \\
\Ye{} improvement over approx21, \msn{80} & 377--651\% & 390--660\% \\
\Ye{} improvement, \msn{151} & 277--390\% & 280--400\% \\
$\bar A$ improvement & 149--291\% / 256--357\% & 150--290\% / 260--360\% \\
\enuc{} improvement & 249--419\% / 284--772\% & 250--450\% / 280--750\% \\
$\nu$-loss crossover & better than approx21 for dt $\le 0.1$~s;
  \textbf{worse} for dt $\in \{1, 10, 100\}$~s & comparable at 0.1~s, worse
  after \\
per-isotope mean $\Delta X$ band (\msn{80}, dt $= 10^{-3}$~s) & 90\% of
  isotopes in $[10^{-4}, 10^{-1}]$, median $3.2 \times 10^{-3}$ & most in that
  band (their Fig.~3) \\
\bottomrule
\end{tabularx}
\endgroup

Two things stand out.

\textbf{The reproduction is exact, not approximate.} The Step-2 pass criterion
was ``within ${\sim}2\times$''; the result was 1.000 in all 108 comparisons and
a bit-identical model handshake. That establishes the baseline as a
\emph{fixed, checkable reference} rather than a quoted number --- this project
can rerun it against any future change.

\textbf{The $\nu$-loss crossover is the baseline's honest weakness}, reported
by the authors and reproduced here: at long timesteps the NNN is \emph{worse}
than approx21 on neutrino losses. Any claim of improvement must be made per-dt,
because the sign of the comparison changes across the grid.

% =====================================================================
\subsection{NuGNN, the closer competitor}
\label{sec:nugnn}

Kim, Chae, Ko, Mumpower \& Smith 2026, \emph{NuGNN: a Graph Neural Network for
Nuclear Reaction Network Equations} (\arxiv{2606.04491}, preprint)
\citep{kim2026} is the nearer prior art: a graph neural network surrogate for
a 690-isotope X-ray-burst network, hence the same architectural family as this
project.

Its documented difficulty is the one that shaped this project's invariant \#4:
attempted enforcement of $\sum\Delta X = 0$ (as a loss term or an output
activation), which requires inverting the signed-log label transform,
\textbf{made training unstable and was dropped}; $\sum\Delta X \approx 0$ is
learned from data (the paper: ``sufficiently well''), backed by a rollout retry
with smaller $\Delta t$ when the predicted sum deviates significantly
(\citealt{kim2026}, \S Training Method and \S Conclusions;
\S\ref{sec:decodestep}). That is a useful competitor to have --- it establishes
that the GNN framing is live and publishable, and it leaves the
\emph{exact-conservation} claim open.

The novelty position this leaves, stated narrowly (and re-checkable at every
phase boundary by the \code{novelty-checker} sweep, which \code{CLAUDE.md}
makes mandatory):

\begingroup\small
\begin{tabularx}{\textwidth}{@{}L{2.1cm} L{2.6cm} L{4.6cm} Y@{}}
\toprule
\textbf{axis} & \textbf{NNN} & \textbf{NuGNN} & \textbf{this project} \\
\midrule
architecture & MLP & GNN, heterogeneous isotope/reaction bipartite, learned
  $(N,Z)$ embeddings & GNN, heterogeneous, raw $(Z,N)$ node features ---
  \emph{not} a differentiator by itself \\
conservation & $\sum X = 1$ by softmax & enforcement through the inverse
  signed-log made training unstable; shipped unenforced (learned + retry) &
  \textbf{exact by construction, all three laws} \\
target & final $X$ & $\Delta X$ (as signed offset-log of $\Delta X/X$) &
  \textbf{per-reaction flux $\Phi$} \\
$\Delta t$ handling & 9 models & one model, $\Delta t$ as input
  ($10^{-15}$--$10^{-4}$~s) & one model, conditioned --- \emph{not} a
  differentiator vs NuGNN \\
reported error & \Ye{} 0.4--0.75\% & MAE 0.037 (offset-log metric); XRB
  rollout deviations mostly $< 15\%$ & --- \\
size transfer & retrain & not addressed & $(Z,N)$ embedding --- a stated
  falsifier \\
\bottomrule
\end{tabularx}
\endgroup

(The NuGNN column was corrected in the 2026-08-18 audit: the earlier table left
its $\Delta t$ handling and node features blank, which overstated this
project's differentiators --- the surviving ones are exact three-law
conservation, the flux-space target, and the size-transfer question.)

% =====================================================================
\subsection{What ``beating it'' has to mean}
\label{sec:beating}

This is the part that needs deciding before Phase 1 rather than after, and it
is not decided anywhere in the repo. The trap: an emulator can be better on the
metric the baseline optimises and worse on the thing that matters, or vice
versa.

A defensible comparison protocol needs at least these five components, and none
of them is currently specified.

\textbf{1. The same test set, and the pathology declared.} The shipped test
sets are labelled by the same bbq that produced the training labels, so they
carry the same pathology above 5~GK (\S\ref{part:pathology}). A comparison in
that band measures agreement with a bug. Either restrict to $\Tn < 5$ or report
the band separately with the caveat attached.

\textbf{2. Per-dt reporting.} The $\nu$-loss crossover
(\S\ref{sec:nnnachieves}) proves that an aggregate over the dt grid can hide a
sign change. Every metric, every dt.

\textbf{3. \Ye{} in absolute units, not relative improvement.} ``377--651\%
improvement over approx21'' is a ratio of losses whose definition is
upstream's. The project's own quantity is $|\dYe|$ per step, absolute, against
the $3 \times 10^{-6}$ gate and the $2.4 \times 10^{-3}$ baseline
(\S\ref{sec:nnncritical}). Both should be reported; only the second is
comparable across papers.

\textbf{4. Rollout, not just one-step.} The NNN's quantitative results are
one-step; the paper reports a 1000-step constant-$(T, \rho)$ rollout only
qualitatively (\Ye{} error stays below approx21's, the other metrics become
comparable or worse) and argues that recursion is not the intended use
\citep[\S 4]{grichener2025}. The scientific use is a \emph{sequence} of steps
($N \approx 1.6 \times 10^3$), where the accumulation character --- systematic
vs random-walk vs contractive --- dominates the outcome (Tier~0 \S IV.3). A
one-step comparison cannot distinguish a model that drifts from one that does
not. \textbf{This is where exact conservation should show its value}, and it
is measurable only in rollout. That the accumulation slope is still unmeasured
(\foreignreg{2}{R-21}) makes this the highest-leverage missing experiment in
the project.

\textbf{5. Trivial baselines, which nobody has run.} \foreignreg{2}{R-16}:
\emph{persistence} ($X' = X$), \emph{nearest neighbour} in the training set,
and above all \textbf{NSE as a predictor}. If NSE predicts the labels well at
$\Tn > 5$ (it does not on this data, because the labels are displaced --- but
it might on physically-correct labels), then the emulator's real domain is
$\Tn < 5$ and its competitor there is a Saha solve, not a neural network.
Running three cheap baselines could redefine the problem. Registered as
\reg{T-08}.

Add one more that follows from \S\ref{sec:verdict}: \textbf{6. Cluster-robust
error bars over trajectories}, so that ``better'' is a claim with an interval
rather than a comparison of two medians.

% =====================================================================
\subsection{The protocol, written out}
\label{sec:protocol}

\S\ref{sec:beating} names six components. Since the point of naming them is to
fix them before Phase 1, here is the concrete version --- a specification that
could be adopted as an ADR today, with the open decisions marked.

\textbf{Test sets.}

\begingroup\small
\begin{tabularx}{\textwidth}{@{}L{3.2cm} L{5.0cm} Y@{}}
\toprule
\textbf{set} & \textbf{source} & \textbf{purpose} \\
\midrule
\textbf{held-out Sobol states} & \code{state\_id} split from the training
  corpus & in-distribution, matches the NNN's own protocol exactly \\
\textbf{relaxed manifold} & the \S\ref{sec:campaign} campaign, trajectories
  held out from any training use & the deployment distribution
  (\S\ref{sec:twodistributions}) --- the NNN was never evaluated here, and this
  is where a conservation-exact model should win \\
\textbf{a real MESA track} & a 20~\Msun{} (and ideally $\ge 3$ progenitors,
  \reg{T-22}) evolution's actual zone states & the \code{CLAUDE.md} Sobol
  $\to$ real gate \\
\textbf{NSE-reachable hot states} & $\Tn \ge 5$ & reported \textbf{separately
  and with the pathology caveat} (\S\ref{part:pathology}), or excluded per
  \reg{T-03} \\
\bottomrule
\end{tabularx}
\endgroup

\textbf{Metrics, all reported per (network, $\Delta t$, \Tn{} stratum):}

\begingroup\small
\begin{tabularx}{\textwidth}{@{}L{4.6cm} Y@{}}
\toprule
\textbf{metric} & \textbf{why} \\
\midrule
$|\dYe|$ \textbf{absolute}, median and p99 & the project's quantity;
  comparable across papers unlike a loss ratio \\
per-isotope $\Delta X$ distribution & comparable to the paper's Fig.~3 band \\
\enuc{} relative error & feeds the energy equation \\
\epsnu{} relative error & \textbf{the NNN's known weak spot} (the crossover,
  \S\ref{sec:nnnachieves}) \\
$\sum X - 1$ and $\min X$ & conservation and positivity; the NNN's softmax
  gives it the first for free, so this is a \emph{tie} not a win --- the win is
  charge \\
$\sum Z\,\dd Y + \dd N_{\mathrm{lep}}$ & \textbf{the charge-to-lepton ledger
  the baseline has no analogue for} \\
\bottomrule
\end{tabularx}
\endgroup

\textbf{Rollout, which is the discriminating experiment.} $N \approx 10^3$
autoregressive steps on the relaxed manifold, reporting the \emph{accumulation
slope} of $|\dYe|$ vs $N$ on a log-log plot. Three outcomes, and the project's
whole conservation argument predicts the first:

\begin{itemize}
\item slope $\approx \tfrac12$ $\to$ random-walk accumulation; the operative
  gate can relax toward $5 \times 10^{-5}$ (\foreignreg{2}{R-21});
\item slope $\approx 1$ $\to$ systematic; the $3 \times 10^{-6}$ gate stands;
\item slope $\approx 0$ $\to$ contractive; the attractor is doing the work, and
  the model's per-step accuracy matters less than anyone thought.
\end{itemize}

\textbf{Baselines}, all of them (\reg{T-08}): the NNN, \code{approx21}
(strictly the 22-isotope \code{approx21\_cr60\_plus\_co56} the NNN paper
compares against, \S 2.3), persistence, nearest-neighbour, and \textbf{NSE} ---
the last on the hot strata, where \S\ref{sec:nsetransition} says it is the real
incumbent for NSE-switching hosts.

\textbf{Uncertainty.} Cluster bootstrap over trajectories
(\S\ref{sec:statspackage}), reported as intervals, on every number.

\textbf{Open decisions this specification does not make}, and which a person
should:

\begin{enumerate}
\item whether $\Tn \ge 5$ is in scope at all (\reg{T-03});
\item whether the headline comparison is against the NNN's \emph{published}
  numbers or a \emph{retrained} NNN on identical splits --- the second is
  fairer and costs a training run;
\item what counts as ``beating'': strictly better on every metric, or better on
  \Ye{} and conservation with parity elsewhere. The second is the honest target
  and should be declared \textbf{before} the numbers exist.
\end{enumerate}

Registered as \reg{T-18}.

% =====================================================================
\subsection{Self-check}
\label{sec:selfcheck-baseline}

\begin{selfcheck}
\item % source VI.6 Q1
  Describe the NNN's architecture and identify the single line that makes it
  conserve mass.
\item % source VI.6 Q2
  Which conservation law does the NNN \emph{not} have, and what is the measured
  consequence in absolute units?
\item % source VI.6 Q3
  The NNN's loss is L1 in log space. Give one strong reason for that choice and
  one structural cost.
\item % source VI.6 Q4
  Why is the $\nu$-loss crossover the most important row in the reproduction
  table for anyone claiming an improvement?
\item % source VI.6 Q5
  Name the five components a defensible ``we beat the baseline'' claim
  requires, and say which one exact conservation is expected to win on.
\item % source VI.6 Q6
  What would running NSE-as-a-predictor as a baseline potentially do to the
  project's scope?
\end{selfcheck}
